In tabular maze worlds, planning pays for itself when the model is right and punishes the agent when it is wrong in the way our setting makes it wrong: a last-outcome model under 0.3 slip makes more planning steps reduce reward until the advantage over Q-learning is gone (at $n\ge50$ the difference from Q-learning is within noise on 10 seeds). Dyna-Q+ clearly improves the response to a stale model and an opened shortcut, but its gain depends on the untried-action initialisation, costs reward under stochasticity, and vanishes or collapses outside a narrow range of the bonus weight. We could not show that Dyna-Q's post-change reward falls with $n$ on the blocking maze, and prioritized sweeping was not distinguishable from Dyna-Q in cumulative or post-change reward at $n=10$ (it did earn more early on). These are results for 20 and 10 seeds on five small mazes with untuned reimplementations; testing whether they carry over to larger tasks and learned stochastic models is the natural next step.
